% =====================================================================
%  EduTwin - Team Beta
%  Resource Recommendation Engine
%  Technical design, architecture, and integration reference
%
%  Compile with:  pdflatex main.tex   (run twice for the table of contents)
%  Only standard TeX Live packages are used.
% =====================================================================

\documentclass[11pt,a4paper]{article}

\usepackage[T1]{fontenc}
\usepackage[utf8]{inputenc}
% Optional: nicer fonts if available, silently skipped on minimal installs.
\IfFileExists{lmodern.sty}{\usepackage{lmodern}}{}
\usepackage[margin=1in]{geometry}
\usepackage{xcolor}
\usepackage{titlesec}
\usepackage{tabularx}
\usepackage{booktabs}
\usepackage{array}
\usepackage{enumitem}
\usepackage{fancyhdr}
\usepackage{listings}
\usepackage{parskip}
\usepackage[hidelinks]{hyperref}

% ---------------------------------------------------------------------
% Colours
% ---------------------------------------------------------------------
\definecolor{navy}{HTML}{1F3864}
\definecolor{slate}{HTML}{44546A}
\definecolor{rule}{HTML}{BFBFBF}
\definecolor{headshade}{HTML}{E8EDF5}
\definecolor{altshade}{HTML}{F5F7FA}
\definecolor{codebg}{HTML}{F4F4F4}

% ---------------------------------------------------------------------
% Headings
% ---------------------------------------------------------------------
\titleformat{\section}
  {\normalfont\Large\bfseries\color{navy}}{\thesection.}{0.6em}{}
  [\vspace{2pt}{\color{rule}\titlerule[0.6pt]}]
\titleformat{\subsection}
  {\normalfont\large\bfseries\color{slate}}{\thesubsection}{0.6em}{}
\titlespacing*{\section}{0pt}{18pt}{8pt}
\titlespacing*{\subsection}{0pt}{14pt}{5pt}

% ---------------------------------------------------------------------
% Header / footer
% ---------------------------------------------------------------------
\pagestyle{fancy}
\fancyhf{}
\renewcommand{\headrulewidth}{0.4pt}
\fancyhead[L]{\footnotesize\color{slate}EduTwin \textperiodcentered\ Resource Recommendation Engine}
\fancyfoot[C]{\footnotesize\color{slate}\thepage}

% ---------------------------------------------------------------------
% Inline code and code blocks
% ---------------------------------------------------------------------
\newcommand{\code}[1]{\texttt{\small #1}}

\lstdefinestyle{plain}{
  basicstyle=\ttfamily\scriptsize,
  backgroundcolor=\color{codebg},
  frame=none,
  breaklines=true,
  columns=fullflexible,
  keepspaces=true,
  xleftmargin=10pt,
  framexleftmargin=10pt,
  framexrightmargin=10pt,
  framextopmargin=5pt,
  framexbottommargin=5pt,
  aboveskip=10pt,
  belowskip=10pt
}
\lstset{style=plain}

% ---------------------------------------------------------------------
% Callout box - no extra packages, just xcolor + minipage
% ---------------------------------------------------------------------
\newsavebox{\calloutbox}
\newenvironment{callout}[1]
  {\par\vspace{8pt}\noindent
   \begin{lrbox}{\calloutbox}
   \begin{minipage}{\dimexpr\textwidth-2\fboxsep-8pt\relax}
   \vspace{4pt}
   {\bfseries\color{navy}#1}\par\vspace{4pt}}
  {\vspace{4pt}\end{minipage}%
   \end{lrbox}%
   \noindent\textcolor{navy}{\rule[-\dp\calloutbox]{3pt}{\dimexpr\ht\calloutbox+\dp\calloutbox\relax}}%
   \colorbox{headshade}{\usebox{\calloutbox}}%
   \par\vspace{10pt}}

% ---------------------------------------------------------------------
% Table helpers
% ---------------------------------------------------------------------
\newcommand{\thead}[1]{\textcolor{white}{\bfseries\small #1}}
\newcommand{\headrow}{\rowcolor{navy}}
\usepackage{colortbl}
\renewcommand{\arraystretch}{1.25}

% Principle block
\newcommand{\principle}[2]{%
  \vspace{8pt}\noindent{\bfseries\color{navy}#1}\par
  \vspace{2pt}\noindent\hspace{6pt}%
  \textcolor{headshade}{\rule{2.5pt}{1em}}\hspace{6pt}%
  \begin{minipage}{\dimexpr\textwidth-20pt\relax}#2\end{minipage}\par\vspace{4pt}}

\setlist[itemize]{leftmargin=1.4em,itemsep=2pt,topsep=4pt}
\setlist[enumerate]{leftmargin=1.6em,itemsep=3pt,topsep=4pt}

\begin{document}

% =====================================================================
% TITLE PAGE
% =====================================================================
\thispagestyle{empty}
\vspace*{5cm}
\begin{center}
{\small\color{slate}\bfseries\spaceskip=0.5em EDUTWIN \textperiodcentered\ TEAM BETA}\\[10pt]
{\Huge\bfseries\color{navy} Resource Recommendation Engine}\\[10pt]
{\color{rule}\rule{0.75\textwidth}{0.6pt}}\\[14pt]
{\large\itshape\color{slate} Technical design, architecture, and integration reference}\\[16pt]
\begin{minipage}{0.78\textwidth}
\centering\small
A deterministic, explainable recommendation engine that consumes a
continuously evolving Digital Twin, reads what the learner is asking for
right now, and remembers every resource it has ever shown them.
\end{minipage}\\[3.5cm]
{\small\itshape\color{slate} Second revision}\\[4pt]
{\small\color{slate} AI Research Internship \textperiodcentered\ Multi-Agent Educational Digital Twin System}
\end{center}
\newpage

% =====================================================================
% CONTENTS
% =====================================================================
\tableofcontents
\newpage

% =====================================================================
\section{Overview}

The Resource Recommendation Engine is one component of the EduTwin multi-agent
educational system. Its job is to recommend educational resources --- videos,
playlists, courses, articles, tutorials, documentation, research papers, and
practice platforms --- that are genuinely useful, verified as reachable, and
personalised to the individual learner.

It is not a search wrapper and it is not a prompt. It is a deterministic
pipeline that reads both the learner's stored profile and what they have just
asked for, retrieves candidates through five progressively relaxed strategies,
verifies them against the live internet, scores them with transparent
arithmetic, refuses to return anything off-topic, and remembers what it has
already shown each person.

\subsection{The research question}

\begin{callout}{Guiding question}
How should AI agents consume a continuously evolving Digital Twin to produce
personalised, explainable, and adaptive educational guidance that outperforms
a traditional stateless LLM assistant?
\end{callout}

\textbf{Core thesis:} personalisation should be an architectural property of
the AI system, not merely a prompt-engineering trick.

\subsection{How this component answers it}

The clearest demonstration is the learner history system. A stateless
assistant has no memory of what it recommended; this engine stores it, keyed
to the learner, in a real database. The difference is directly observable:

\vspace{4pt}
\begin{tabularx}{\textwidth}{>{\raggedright\arraybackslash}p{3.6cm} >{\raggedright\arraybackslash}X >{\raggedright\arraybackslash}X}
\headrow \thead{Learner action} & \thead{Stateless LLM assistant} & \thead{This engine} \\
\midrule
Asks the same question twice & May give different answers at random & Identical answer, by design \\
\rowcolor{altshade} Says ``recommend something else'' & Often repeats itself --- no memory of what it said & Excludes everything shown; searches with varied queries for genuinely new material \\
Rejects a resource & Forgotten immediately & Permanently excluded, across sessions \\
\rowcolor{altshade} Asks why a resource was chosen & Post-hoc rationalisation & Exact score breakdown, generated from the arithmetic \\
\bottomrule
\end{tabularx}
\vspace{8pt}

\textbf{Reproducibility is deliberate.} The same learner asking the same
question receives the same answer, so any recommendation can be explained,
audited, and defended. Variety is available --- but it must be explicitly
requested, not delivered by accident.

\begin{callout}{Stated honestly}
The learner history system requires a learner identifier, and the integrated
system does not currently supply one. In the live pipeline, this feature is
inert. Section 10 covers this in full.
\end{callout}

% =====================================================================
\newpage
\section{The governing architectural principle}

\begin{callout}{Non-negotiable}
Large language models must NOT make the final recommendation decision.
\par\vspace{4pt}
LLMs are used only for: understanding natural-language requests, extracting
structured fields from prose, interpreting unstructured resource content, and
generating explanations. All filtering, scoring, ranking, and the final
decision are 100\% deterministic code.
\end{callout}

This constraint drives every design decision in the component. Scoring is
plain arithmetic rather than ``ask a model which is best,'' because arithmetic
is reproducible, debuggable, and can be shown to a learner as a justification.

\subsection{Where the principle proved its worth}

During development a significant bug was found: every resource discovered from
the internet was claiming to be a video and collecting thirty unearned scoring
points. The cause was a single line assuming that searching for videos
guarantees video results --- in reality the search returned a YouTube
playlist, a social-media discovery page, and two documentation pages, all
labelled ``video.''

The bug was visible because the score breakdown is printed transparently: a
resource showing \code{format 30} beside a category-listing URL is obviously
wrong. A system that asked a model to rank resources would have produced a
confident, plausible, and equally incorrect answer --- with nothing to
inspect.

\subsection{Division of responsibility}

\vspace{4pt}
\begin{tabularx}{\textwidth}{>{\raggedright\arraybackslash}p{5.2cm} >{\raggedright\arraybackslash}p{4.2cm} >{\raggedright\arraybackslash}X}
\headrow \thead{Concern} & \thead{Handled by} & \thead{Why} \\
\midrule
Understanding what the learner said & LLM (Orchestrator) & Natural language is exactly what models are good at \\
\rowcolor{altshade} Extracting fields from a prose briefing & LLM (context extractor) & Same reason; heavily guard-railed \\
Reading the learner's immediate request & LLM (request extractor) & Same reason; falls back to deterministic keywords \\
\rowcolor{altshade} Interpreting video and article content & LLM (content analyzers) & Unstructured content requires comprehension \\
Inferring difficulty and format & Deterministic code & Feeds sixty of one hundred scoring points \\
\rowcolor{altshade} Judging topic relevance & Deterministic code (embeddings) & Decides whether a resource is shown at all \\
Filtering, scoring, ranking & Deterministic code & Must be reproducible and explainable \\
\rowcolor{altshade} The final recommendation & Deterministic code & The core claim of the research \\
\bottomrule
\end{tabularx}

% =====================================================================
\newpage
\section{How the engine works}

A single function call runs the entire pipeline. The caller supplies two
strings --- the learner's stored briefing and the message they just typed ---
and receives ranked, explained, verified recommendations.

\subsection{The pipeline}

\begin{lstlisting}
Two strings arrive: the briefing, and the learner's message
  |  extract learner state from the briefing   (cache -> Gemini)
  |  extract the request from the message      (cache -> Gemini -> keywords)
  |  APPLY REQUEST: the request wins, memory fills the gaps
  |  normalise (field aliases, reject invalid values)
  |  load this learner's history from PostgreSQL
  |  build the search requirement (topic, level, format)
  |
  +- TIER 1  exact topic + exact level           | each filtered against
  +- TIER 2  exact topic, any level              | history; stops at the
  +- TIER 3  semantic match (neural)             | first tier that yields
  +- TIER 4  entire catalogue - ONLY if no topic | unseen resources
  +- TIER 5  external internet discovery         |
        +- query variation (8 templates, neutral first)
        +- analyse   (infer difficulty and format)
        +- verify    (structure -> blocklist -> homepage -> parallel HTTP)
        +- store     (save all verified to the shared catalogue)
  |  score every resource  (format 30 / level 30 / duration 20 / goal 20)
  |  TOPIC GATE: off-topic resources score zero
  |  rank, drop anything below 50, cap at 5 or the number requested
  |  record what was shown
  |  deep content analysis of the top resource only (cached, shared)
  |  format as readable text
\end{lstlisting}

\subsection{Tiered retrieval}

Each tier is a deliberate relaxation of constraints. The engine never returns
nothing when something reasonable exists --- but it never pretends a loose
match is a tight one either.

\textbf{Tier 4 only runs when no topic was specified.} A catch-all is right
when the engine has no idea what the learner wants; it is wrong when they
named one. A real run returned all forty-seven machine learning resources to a
learner asking for calculus, because the catch-all matched everything and
internet discovery therefore never ran.

Every tier is also filtered against the learner's history. That is what makes
an exhausted learner fall through to fresh discovery rather than receiving an
empty result or a repeat.

\subsection{Query variation}

When a learner has seen everything the obvious search returns, the engine
escalates through eight fixed search phrasings:

\begin{lstlisting}
1.  {topic} {level} {format}       5.  how to {topic} {format}
2.  {topic} {format}               6.  {topic} full guide {format}
3.  best {topic} {format}          7.  learn {topic} from scratch {format}
4.  {topic} explained {format}     8.  {topic} crash course {format}
\end{lstlisting}

\textbf{The neutral phrasings come first, deliberately.} An earlier version
opened with the learning-oriented templates, which produced ``learn messi from
scratch video'' for a non-academic topic --- unnatural phrasing that searches
badly. Templates one to four work for any subject; five to eight are the
academic escalation, used only when the plain searches run dry.

In the normal case only the first template runs, so the mechanism costs
nothing until it is needed. Because the list is fixed and ordered --- not
generated by a model --- the same topic with the same exclusion set always
produces the same query sequence. Adaptive behaviour without sacrificing
reproducibility.

% =====================================================================
\newpage
\section{Reading the learner's request}

A learner's goal is long-term: ``become a machine learning engineer.'' Their
request is immediate: ``show me videos about calculus.'' An early version of
the engine read only the goal, so a learner asking about calculus received
machine learning resources.

\begin{callout}{The rule}
THE REQUEST WINS, MEMORY FILLS THE GAPS.
\par\vspace{4pt}
Each field is taken from the learner's message when it states one, and from
the Digital Twin when it does not. ``Recommend videos about calculus''
searches calculus at the learner's stored level and duration. ``Recommend
something'' uses the stored profile entirely.
\end{callout}

When the request overrides the stored goal, the response says so rather than
silently searching for something the profile does not expect:

\begin{lstlisting}
Searching for 'calculus' as requested, rather than the stored goal
'machine learning'.
\end{lstlisting}

\subsection{Three extraction layers}

Reading the request is attempted three ways, cheapest and most reliable first:

\vspace{4pt}
\begin{tabularx}{\textwidth}{>{\raggedright\arraybackslash}p{2.4cm} >{\raggedright\arraybackslash}X >{\raggedright\arraybackslash}X}
\headrow \thead{Layer} & \thead{What it does} & \thead{When it runs} \\
\midrule
1 \textperiodcentered\ Cache & Returns a previously extracted result for this exact message & Always tried first --- free, instant, immune to API outages \\
\rowcolor{altshade} 2 \textperiodcentered\ Gemini & Full understanding: any phrasing, any language, expands abbreviations & When the cache misses \\
3 \textperiodcentered\ Keywords & Deterministic pattern matching on common request shapes & Only when Gemini is unavailable \\
\bottomrule
\end{tabularx}
\vspace{8pt}

\textbf{The third layer exists because the API failed repeatedly in real use}
--- quota exhausted, servers overloaded, and once an unset key. With no
extraction at all, the learner state is empty and every recommendation scores
zero out of one hundred. Keyword results are deliberately not cached, so a
crude answer is never reused in place of the better one once the API recovers.

\subsection{Abbreviations and counts}

The extracted topic is used directly as an internet search query, so shorthand
is expanded: ``calc'' becomes ``calculus'', ``ml'' becomes ``machine
learning'', ``linalg'' becomes ``linear algebra''. Proper nouns --- names of
people, games, products --- are left exactly as written.

A learner may also ask for a specific number. ``Give me ten videos about
python'' returns up to ten rather than the default five. The value is clamped
between one and twenty, so an extreme request cannot trigger hundreds of
network verifications.

\textbf{Asking for ten does not guarantee ten.} Only resources that pass the
topic gate and clear the quality floor are returned; padding the list with
weak matches would defeat the filter. The response says so plainly: ``Only six
strong matches were found, fewer than the ten requested.''

% =====================================================================
\newpage
\section{Scoring, the topic gate, and the quality floor}

One hundred points, four factors, plain arithmetic:

\vspace{4pt}
\begin{tabularx}{\textwidth}{>{\raggedright\arraybackslash}p{4cm} >{\centering\arraybackslash}p{1.8cm} >{\raggedright\arraybackslash}X}
\headrow \thead{Factor} & \thead{Points} & \thead{Awarded when} \\
\midrule
Format match & 30 & The learner's preferred format is known and matches \\
\rowcolor{altshade} Level match & 30 & The learner's level is known and matches \\
Duration match & 20 & The learner's preferred duration is known and matches \\
\rowcolor{altshade} Goal relevance & 20 & The resource is genuinely about the requested topic \\
\textbf{Total} & \textbf{100} & \\
\bottomrule
\end{tabularx}
\vspace{8pt}

A factor earns points only if the learner's value is known \emph{and} matches.
It never guesses a match for missing data. A separate confidence value between
zero and one reports how much of the learner's profile was actually available,
so a low score caused by an unknown profile is distinguishable from a low
score caused by a genuinely poor match.

\subsection{The topic gate}

\begin{callout}{The rule}
When the learner has named a topic, a resource that is not about that topic
scores ZERO --- whatever else it matches.
\end{callout}

A machine learning video is not a useful answer to ``recommend videos about
calculus'' just because the format happens to be right. Before this gate
existed, a real run returned exactly that, at thirty out of one hundred,
purely on the format match.

\textbf{Relevance is judged semantically, not by string equality.} A resource
stored under ``machine learning'' still counts for a request about ``machine
learning basics''. Exact matching would reject it, throw away almost the
entire catalogue, and force a web search on nearly every request.

The gate is applied identically in the scorer and in the decision-record
builder, so a record's score and its stated reasons can never disagree. An
off-topic resource's reason reads: ``Not about `calculus', which is what was
requested.''

\subsection{The quality floor}

Resources scoring below fifty are not shown. A weak match is worse than an
honest ``nothing good found'' --- it wastes the learner's time and makes the
engine look careless.

With one safeguard: \textbf{if nothing clears the bar, the best available are
returned anyway, with a message saying none scored fifty or above.} Never
return nothing when something exists, but never present a weak match as though
it were a good one either.

% =====================================================================
\newpage
\section{Quality, safety, and verification}

A recommendation engine that surfaces dead links, pirated material, or content
farms is worse than no engine at all. Four gates run on every candidate,
ordered cheapest first so that a network request is never spent on a resource
that was always going to be rejected.

\vspace{4pt}
\begin{tabularx}{\textwidth}{>{\raggedright\arraybackslash}p{3.4cm} >{\raggedright\arraybackslash}X >{\raggedright\arraybackslash}p{3.2cm}}
\headrow \thead{Gate} & \thead{Check} & \thead{Cost} \\
\midrule
1 \textperiodcentered\ Structural & Title, URL, and topic must all be present & Free \\
\rowcolor{altshade} 2 \textperiodcentered\ Domain policy & URL must not match the blocklist & Free \\
3 \textperiodcentered\ Homepage & A site's front door is not a resource & Free \\
\rowcolor{altshade} 4 \textperiodcentered\ Reachability & Real HTTP request; status code below 400 & One network call \\
\bottomrule
\end{tabularx}
\vspace{8pt}

\subsection{The homepage gate}

Added after a real run recommended a bare site homepage as a machine learning
resource scoring fifty out of one hundred, \textbf{purely because the page
title contained the word ``learning''.} A homepage has no specific content, no
difficulty and no single topic. There is nothing to recommend, whatever it
scores.

\subsection{Domain policy}

Real discovery runs surfaced pirated-course sites, social-media discovery
pages, anonymous free-hosting subdomains, a video-scraper site, and
stock-media libraries selling machine learning footage --- alongside genuinely
excellent material from Google, Microsoft Learn, and freeCodeCamp. The scoring
system cannot tell them apart, because quality is not one of its four factors.

The policy holds two lists with deliberately different powers:

\begin{itemize}
\item \textbf{Blocked domains} --- piracy, social feeds, anonymous hosting,
content farms, stock media, and vendor product pages. Refused outright.
\item \textbf{Trusted domains} --- universities, official documentation,
established platforms. Marked as trusted, but nothing is rejected for failing
the check.
\end{itemize}

\textbf{A domain on neither list is allowed through as neutral.} Blocking is
deliberate, trust is a bonus, and everything else is permitted --- an
allowlist-only approach would starve the catalogue of good material from sites
nobody thought to list.

This list needs ongoing curation. New categories of junk kept appearing
throughout development, each discovered by inspecting real output rather than
anticipated in advance.

\subsection{Parallel verification}

Verifying URLs one at a time means adding up the waiting. Fifteen URLs with a
five-second timeout each produced waits of up to two minutes in real runs,
since a dead link burns the full timeout before failing.

The network gate now runs across a thread pool. Threads are the correct tool
here specifically because this is network waiting rather than calculation ---
Python releases its global lock while waiting on a socket, so the other
threads genuinely run. Result ordering is preserved, which matters because
ranking depends on it.

\begin{callout}{Measured result}
8 URLs, sequential: 5.4 seconds. Parallel: 0.8 seconds. 7.1$\times$ faster,
with identical results and preserved ordering.
\end{callout}

\subsection{Ongoing catalogue audit}

Links die. A separate audit script re-checks stored resources on a schedule
--- never during a learner's request. A resource is deleted only after five
consecutive failures; the count resets to zero on any success, so a link that
usually works never accumulates strikes.

\begin{callout}{Safety net}
If more than 70\% of a batch fails at once, that is almost certainly a local
network problem rather than that many simultaneously dead links. The run
aborts and records nothing.
\par\vspace{4pt}
This guard was written after two total network outages during development.
Without it, either one would have deleted the entire catalogue.
\end{callout}

% =====================================================================
\newpage
\section{Data design}

Four tables in a live cloud PostgreSQL database. Each carries constraints
chosen to make invalid states structurally impossible rather than something
the application layer must remember to prevent.

\vspace{4pt}
\begin{tabularx}{\textwidth}{>{\raggedright\arraybackslash}p{4.4cm} >{\raggedright\arraybackslash}X >{\raggedright\arraybackslash}p{4cm}}
\headrow \thead{Table} & \thead{Holds} & \thead{Shared across learners?} \\
\midrule
\code{resources} & Every resource the engine knows about & Yes --- the catalogue grows with use \\
\rowcolor{altshade} \code{recommendation\_history} & What each learner has been shown or has rejected & No --- keyed per learner \\
\code{resource\_analysis} & Cached deep content analysis & Yes --- analysed once, ever \\
\rowcolor{altshade} \code{extraction\_cache} & Cached LLM extractions, keyed by input text & Yes \\
\bottomrule
\end{tabularx}
\vspace{8pt}

\subsection{recommendation\_history --- learner memory}

The table that makes the engine stateful, and therefore the table that carries
the research claim. Four design decisions are worth explaining:

\vspace{4pt}
\begin{tabularx}{\textwidth}{>{\raggedright\arraybackslash}p{5cm} >{\raggedright\arraybackslash}X}
\headrow \thead{Decision} & \thead{Rationale} \\
\midrule
Composite primary key on (twin\_id, resource\_url) & One row per learner per resource, ever. A learner cannot simultaneously be ``shown'' and ``rejected'' for the same resource --- rejecting updates the existing row. Contradictory state is impossible rather than something a query must resolve. \\
\rowcolor{altshade} Foreign key targets url, not the integer id & Serial IDs are reassigned whenever the table is rebuilt, which would silently repoint a learner's history at the wrong resources. URLs are stable identities. \\
event is text with a CHECK, not a boolean & Three states are needed: never shown (no row), shown (mild signal), rejected (permanent). A boolean cannot express that. The constraint also stops typos silently never matching a query. \\
\rowcolor{altshade} ON DELETE CASCADE & Deleting a resource removes its history rows, keeping the data consistent. The trade-off is documented in the limitations section. \\
\bottomrule
\end{tabularx}
\vspace{8pt}

\subsection{The two caches}

Deep content analysis costs an API call per resource; so does reading a
briefing or a request. Both are cached, and both caches are shared across
every learner --- work paid for by one benefits all.

\textbf{Only successful results are cached.} Storing a failure would freeze a
temporary API outage permanently, and that resource could never be analysed
again. This distinction proved its value immediately: a live ``service
unavailable'' error was served correctly from cache moments later, with no API
call at all.

The extraction cache is keyed by a hash of the input text rather than the text
itself --- a briefing can run to thousands of characters, far too long for a
primary key. Beyond saving API calls, it closes a reproducibility gap: the
same briefing now always yields the same extraction, which prose extraction
alone could not guarantee.

% =====================================================================
\newpage
\section{Integration contract}

The engine exposes one simple entry point. Everything behind it is internal
and free to change.

\begin{lstlisting}
from orchestrator_interface import recommend_text

text = recommend_text(briefing_string, user_request_string)
\end{lstlisting}

\subsection{Two inputs, both strings}

\vspace{4pt}
\begin{tabularx}{\textwidth}{>{\raggedright\arraybackslash}p{3cm} >{\raggedright\arraybackslash}p{4.6cm} >{\raggedright\arraybackslash}X}
\headrow \thead{Input} & \thead{What it is} & \thead{Why it matters} \\
\midrule
\code{briefing} & The context builder's prose output & The learner's stored goal, level, preferred format and duration \\
\rowcolor{altshade} \code{user\_request} & What the learner just typed & What they want right now, which overrides the stored profile \\
\bottomrule
\end{tabularx}
\vspace{8pt}

Two optional flags: one to exclude everything the learner has already been
shown, and one to skip deep content analysis for a faster response. The output
is ready-to-display text.

\subsection{The structured variant}

A second function returns the same result as a dictionary rather than text.
Use it when the interface needs individual fields --- clickable links, score
badges, a reject button. Those values cannot be recovered reliably from
formatted text.

The response carries the recommendation list, each with its URL, format,
score, full score breakdown, human-readable reasons, and confidence value;
plus the deep analysis of the top resource and a field naming which resource
that analysis belongs to.

\textbf{That last field matters.} The interface should match the analysis to a
resource by URL rather than assuming it belongs to the first item --- if
ranking behaviour ever changes, matching by URL keeps the display correct.

\subsection{Three intents}

\vspace{4pt}
\begin{tabularx}{\textwidth}{>{\raggedright\arraybackslash}p{5.4cm} >{\raggedright\arraybackslash}X}
\headrow \thead{Intent} & \thead{Purpose} \\
\midrule
\code{resource\_recommendation} & Retrieve, score, and rank resources. Returns deep content analysis of the top result. \\
\rowcolor{altshade} \code{reject\_resources} & Record that a learner does not want specific resources. They are never recommended to that learner again. \\
\code{analyze\_resource\_content} & Deep analysis of one specific resource. Built for a per-resource ``What's inside?'' button; not currently used by the interface, but tested and ready. \\
\bottomrule
\end{tabularx}
\vspace{8pt}

\subsection{Requirements of the calling system}

\begin{enumerate}
\item \textbf{Send a learner identifier.} Without one the engine cannot
remember anything --- no history, no exclusions, and rejections fail outright.
See Section 10.

\item \textbf{Decide with the model, not a keyword list.} When a learner asks
for different resources --- in any phrasing, in any language --- the
Orchestrator's language model should recognise that and set the flag. This
engine performs no language understanding on the user's intent, by design.

\item \textbf{Show a loading state.} A first request on a new topic takes
thirty to sixty seconds while the engine searches and verifies the live
internet. Later requests on the same topic are answered from the catalogue and
are fast.

\item \textbf{Expect determinism.} The same learner asking the same question
receives the same answer. Variety requires an explicit flag.
\end{enumerate}

\subsection{Output formatting}

The text output is deliberately compact. An early version printed every
chapter and section with its full paragraph summary --- an eight-section
article buried the recommendations under a wall of text. The current format
shows headings only, capped at five items, and for articles leads with the
sections flagged as relevant to the learner's topic.

% =====================================================================
\newpage
\section{Resilience}

Every external dependency has a fallback that degrades honestly and explains
itself. This was not a theoretical exercise --- the design was validated
repeatedly by genuine infrastructure failures during development.

\vspace{4pt}
\begin{tabularx}{\textwidth}{>{\raggedright\arraybackslash}p{3.4cm} >{\raggedright\arraybackslash}p{4.4cm} >{\raggedright\arraybackslash}X}
\headrow \thead{Dependency} & \thead{When it fails} & \thead{What happens} \\
\midrule
Neural embeddings & Model cannot load & Falls back to bag-of-words similarity; prints which is active \\
\rowcolor{altshade} PostgreSQL & Network, credentials, paused instance & Falls back to an in-memory catalogue; prints the real error \\
Search engine & Blocked, unreachable, rate-limited & Returns no candidates; the request completes with an honest message \\
\rowcolor{altshade} Gemini --- briefing & Outage, quota, unset key & Proceeds with unknown learner state and says so \\
Gemini --- request & Outage, quota, unset key & Falls back to deterministic keyword matching \\
\rowcolor{altshade} Either cache & Read or write fails & Proceeds without the cache \\
\bottomrule
\end{tabularx}
\vspace{8pt}

\subsection{Layered, not merely present}

The request extractor is the clearest example: cache first, then the API, then
deterministic code. Each layer is cheaper and more reliable than the one
before it, and each is genuinely useful on its own.

\begin{callout}{Validated under real failure}
During development the network became entirely unavailable twice --- once
through TLS interception, once through blocked outbound database traffic. The
Gemini API failed twice more, once from an exhausted daily quota and once from
server overload.
\par\vspace{4pt}
In every case the engine continued to return recommendations, printed the
precise cause of each degradation, and did not crash. Every fallback path was
exercised in real conditions, unplanned.
\end{callout}

\textbf{The pattern throughout is the same:} an honest status message
describing what could not be done, rather than a plausible-looking result that
quietly hides a failure. The same discipline governs the ``not yet supported''
responses for book-format resources, for playlists that have no single
transcript, and for pages whose text cannot be extracted.

% =====================================================================
\newpage
\section{The blocking integration issue}

\begin{callout}{A learner identifier is never supplied}
Every test through the integrated system extracts a null learner identifier.
The briefing contains a name, university, field of study and education stage
--- but nothing that identifies the learner to this engine.
\end{callout}

\textbf{The consequence is that the entire learner memory system does nothing
in the live system.}

\begin{itemize}
\item Already-seen resources are not excluded.
\item Rejections fail outright.
\item ``Recommend something else'' returns the same results.
\item Nothing is recorded between sessions.
\end{itemize}

This is the feature that most directly demonstrates the research question ---
the difference between this engine and a stateless assistant --- and in the
integrated system it is inert.

\textbf{It cannot be fixed from inside this component.} The engine already
accepts four different field names for the identifier and issues a warning
when none is present. The calling system must include it.

Everything else in this document describes work that is complete and tested.
This one item is a conversation, not a code change, and it is the single
highest-priority action.

% =====================================================================
\newpage
\section{Known limitations}

Stated plainly. Several are deliberate trade-offs rather than defects, but all
are known and none are hidden.

\subsection{Affecting recommendation quality}

\begin{itemize}
\item \textbf{Duration is never inferred.} Nothing determines whether a
resource is short, medium, or long, so every resource forfeits twenty of one
hundred points and realistic scores cap at eighty.

\item \textbf{Quality is not a scoring factor.} Among resources that pass the
topic gate, nothing distinguishes an excellent resource from a mediocre one
--- no view counts, no reputation, no ratings. The trusted-domain list exists
and is not yet wired into scoring; doing so is the most obvious available
improvement.

\item \textbf{Non-educational topics rank poorly.} A request for sport or
gaming content extracts correctly, but most such content lives on blocked
social domains, and with no quality signal every result scores low. The engine
is tuned for educational topics; this is a scope boundary rather than a
defect.

\item \textbf{Score ties are broken arbitrarily.} The sort is stable, so the
winner among equal scores is whichever was stored first --- not merit.
\end{itemize}

\subsection{Affecting reproducibility}

\begin{itemize}
\item \textbf{Prose extraction is not fully reproducible.} The same briefing
may occasionally yield a different level or format. The recommendation
decision remains fully deterministic given a learner state; it is the
extraction step that varies. The cache removes this for inputs already seen.

\item \textbf{Goal selection is a model judgment.} When the learner's profile
holds several goals, the extractor picks the one presented as current. If the
briefing exposed which goal is active, this could be deterministic instead.

\item \textbf{Degraded mode is non-deterministic.} When the database is
unreachable, discovery adds to the in-memory catalogue mid-session, so later
requests see a different catalogue than earlier ones.
\end{itemize}

\subsection{Affecting data and maintenance}

\begin{itemize}
\item \textbf{The blocklist needs ongoing curation.} New categories of
low-quality domain kept appearing throughout development, each found by
inspecting real output rather than anticipated.

\item \textbf{The blocklist uses substring matching.} Sufficient for filtering
search results, insufficient against a determined actor. Proper URL host
parsing would be the robust fix.

\item \textbf{Cached analyses never expire.} If a video is re-uploaded or an
article rewritten, stale content is served indefinitely. The timestamp needed
for a staleness policy is stored, but no policy uses it.

\item \textbf{Deleting a resource erases learner history for it,} including
rejections. If it is later rediscovered, a learner who rejected it could see
it again.

\item \textbf{The keyword fallback is crude by design.} It handles common
request shapes but would fail on a long conversational message. It only runs
when the API is unavailable.
\end{itemize}

\subsection{Out of scope by decision}

\textbf{Books are deliberately unsupported for content analysis.} Most books
are copyrighted or paywalled with no legitimate full-text access. Rather than
build something that appears to work but does not, book-format resources
receive an explicit ``not yet supported'' message. This was a conscious
decision, not an oversight.

% =====================================================================
\newpage
\section{Development history}

A record of genuine failures encountered and resolved. Kept because the
reasoning is often more instructive than the fix, and because several
represent traps that would otherwise be fallen into again.

\subsection{Correctness}

\vspace{4pt}
\begin{tabularx}{\textwidth}{>{\raggedright\arraybackslash}p{4cm} >{\raggedright\arraybackslash}X}
\headrow \thead{Problem} & \thead{Cause and resolution} \\
\midrule
Unearned scoring points & Discovery labelled every result with the format that had been searched for. A playlist, a social-media page and two documentation pages all claimed to be videos and collected thirty points each. Fixed by leaving format unknown at discovery and inferring conservatively from the URL, including path inspection to distinguish a single video from a playlist. \\
\rowcolor{altshade} Off-topic resources returned & A calculus request returned all forty-seven machine learning resources. Two causes: the catch-all tier matched everything before internet discovery could run, and format-only matches scored thirty with no topic requirement. Fixed by restricting the catch-all and adding the topic gate. \\
A site homepage recommended & A bare homepage scored fifty because its title contained a matching word. Fixed with a homepage gate in the verifier. \\
\rowcolor{altshade} Subject names destroyed & The keyword fallback stripped ``learning'' as a filler word, turning ``machine learning'' into ``machine'' --- and the engine recommended a sewing machine tutorial. Fixed by removing those words from the filler list. \\
A filler word became the topic & ``Recommend something for me'' produced the topic ``for'', and the engine searched the web for it. Fixed by excluding marker words from the extracted topic. \\
\rowcolor{altshade} Twenty-six recommendations returned & The first tier returns everything matching, and the catalogue grows every run. Worse, all twenty-six were recorded as shown, wrongly excluding twenty-one good resources from every future request. Fixed by capping the list before recording history. \\
Verified work discarded & Surplus resources that had each cost a real network request were thrown away, forcing rediscovery later. Fixed by saving everything verified and returning only what the request needs. \\
\bottomrule
\end{tabularx}

\subsection{Environment and tooling}

\vspace{4pt}
\begin{tabularx}{\textwidth}{>{\raggedright\arraybackslash}p{4cm} >{\raggedright\arraybackslash}X}
\headrow \thead{Problem} & \thead{Cause and resolution} \\
\midrule
Silent loss of neural embeddings & A bad import inside a function body fired only when the function was called, producing a misleading ``module not found'' message and a silent fallback to the weaker similarity method for days. Found by listing every import in the project. \\
\rowcolor{altshade} Files documented but absent & Several functions and one entire module were described as complete but were not on disk. Documentation had drifted from reality; the filesystem is the authority. \\
Three distinct network failures & On the same connection: TLS interception broke all HTTPS; the database host resolved to an unroutable address family; and outbound database traffic was blocked entirely. None were code faults. Certificate verification was deliberately not disabled --- for a system whose purpose is verifying resources, that would be the wrong trade. \\
\rowcolor{altshade} API quota and outages & The free tier allows twenty requests per day per model, and each recommendation makes up to four calls --- roughly five test runs. The service also returned overload errors independently. This directly motivated the extraction cache and the keyword fallback. \\
Environment variables lost repeatedly & Session-scoped variables vanish when a terminal closes, and an unset database host silently defaulted to the local machine, making a configuration problem look like an outage. Two further traps: putting the project reference in the database name rather than the username, and an invisible newline stored inside a value. \\
\bottomrule
\end{tabularx}
\vspace{8pt}

\textbf{A recurring theme:} the most costly bugs were not crashes. They were
silent --- a resource collecting points it had not earned, a subject name
quietly truncated, a documented file that did not exist, a configuration value
with an invisible character. Each was found by inspecting output that the
system had been built to expose.

% =====================================================================
\newpage
\section{Outstanding work}

\subsection{Immediate}

\begin{enumerate}
\item \textbf{Resolve the learner identifier} with the Orchestrator developer.
Nothing else matters as much --- see Section 10.
\item \textbf{Verify end-to-end persistence} on an unrestricted network:
discover, save, restart, and confirm the resources are still there.
\item \textbf{Create a shared dependency list} for the whole project. Packages
are currently installed one at a time as errors appear, and every team member
hits the same wall.
\end{enumerate}

\subsection{Short-term improvements, in value order}

\begin{enumerate}
\item \textbf{Wire the trusted-domain list into scoring} as a quality signal
and a tiebreaker between the many resources that currently score identically.
The single biggest available improvement to recommendation quality.
\item \textbf{Infer duration} from video length or estimated reading time ---
recovers twenty points currently forfeited on every resource.
\item \textbf{Continue curating the blocklist} as new categories of
low-quality domain appear.
\end{enumerate}

\subsection{Discussed but not built}

\begin{enumerate}
\item \textbf{Blended recommendations.} After the requested topic, add two or
three resources relevant to the learner's stored goal, clearly labelled. The
ambitious version --- recommending linear algebra because it supports machine
learning --- needs prerequisite knowledge the engine does not have. The
simpler version needs no new knowledge and would work today.
\item \textbf{An intermittent interface failure} was reported but never
diagnosed, because the failure output was never captured. Most likely a
timeout on the slow first request, but this is unconfirmed.
\end{enumerate}

\subsection{Longer-term, dependent on other teams}

\begin{enumerate}
\item \textbf{Real Digital Twin integration.} The learner model stores goals
and preferences as collections of objects; this engine expects single values.
The bridging responsibility must be agreed before integration rather than
discovered during it.
\item \textbf{Knowledge graph integration.} Would allow retrieval to broaden
into prerequisite and adjacent concepts rather than falling straight through
to a web search. Currently blocked: the graph is not persisted between runs
and mastery values are never populated.
\item \textbf{Emit rejection signals upstream.} When a learner rejects
resources, that is real evidence about their preferences. The engine's history
table drives filtering; a memory record could carry the signal to the learner
model, which might conclude the format or level is wrong.
\end{enumerate}

% =====================================================================
\newpage
\section{Design principles}

Extracted from decisions that proved correct under real use. If this component
is continued or reimplemented, these are the rules that kept it honest.

\principle{Never claim what you have not verified}{The format bug awarded
thirty points for an assumption. Every field is now either known and correct,
or explicitly unknown. Unknown scores zero, which is the honest outcome ---
and a resource scoring fifty truthfully is more useful than one scoring eighty
on a guess.}

\principle{A wrong answer is worse than no answer}{The topic gate and the
quality floor both exist because returning something plausible but irrelevant
wastes the learner's time and makes the engine look careless. But never return
nothing when something exists --- say plainly that the matches are weak.}

\principle{Make invalid states structurally impossible}{A composite primary
key means a learner cannot be both shown and rejected the same resource. A
uniqueness constraint means duplicates cannot accumulate. Rules enforced by
the database survive buggy code, other people's scripts, and manual edits.}

\principle{Degrade honestly, never crash}{Every external dependency has a
fallback that explains itself. Validated repeatedly by real network, database
and API failures --- the engine kept working and said precisely what was
wrong.}

\principle{Layer the fallbacks}{Cache, then API, then deterministic code. Each
layer is cheaper and more reliable than the one before it, and each is
genuinely useful on its own rather than being a token gesture.}

\principle{Cheapest checks first}{Structural, blocklist and homepage checks
run before any network request. Cache lookups run before any API call. Work
that can be avoided should be avoided before it is paid for.}

\principle{Pay once, share the result}{Verified resources and cached analyses
are stored in shared tables. Work done for one learner benefits every learner,
and the catalogue grows more valuable with use.}

\principle{Reproducibility is a feature, not a constraint}{The same input
produces the same output, so a recommendation can be explained, audited, and
defended. Variety is available on request --- never by accident.}

\principle{Documentation drifts; the filesystem does not}{Several files this
project's own notes described as complete were not on disk. Verify against
reality before building on an assumption.}

\principle{Say when you do not know}{Several bugs were diagnosed only after
asking for the actual output rather than reasoning from a description of it. A
confident guess is worse than an admitted gap.}

\end{document}